\documentclass[12pt,a4paper]{report}

\usepackage[utf8]{inputenc}
\usepackage[T1]{fontenc}
\usepackage{lmodern}
\usepackage[margin=1in]{geometry}
\usepackage[draft]{graphicx}
\usepackage{hyperref}
\usepackage{enumitem}
\usepackage{longtable}
\usepackage{array}

\hypersetup{
  colorlinks=true,
  linkcolor=black,
  urlcolor=blue,
  pdftitle={CourShera Final Report},
  pdfauthor={B2 Group 3}
}

\title{CourShera Final Report\\CSE 326: Information System Design Sessional}
\author{B2 Group 3}
\date{April 2026}

\begin{document}

\maketitle
\tableofcontents
\clearpage

% =========================================================
\chapter{Introduction}

\textbf{Presented by:} [Member Name, Student ID]

\section{Project Overview}
CourShera is a web-based learning platform inspired by Coursera, developed as part of the CSE 326 ISD Sessional project. The system enables learners to browse courses by category, view course details, register/login (local and Google OAuth), manage profile information, enroll in courses through payment verification, and access quiz-related learning activities.

The goal of the project is to design and implement an end-to-end educational platform that demonstrates complete software engineering workflow: requirements analysis, modeling, UI design, API design, implementation, and deployment readiness.

\section{Problem Statement}
Students and self-learners need a single, organized platform where they can:
\begin{itemize}[leftmargin=*]
  \item Discover relevant courses quickly.
  \item Track progress in enrolled courses.
  \item Complete secure enrollment/payment workflows.
  \item Manage their learning profile and credentials.
\end{itemize}

Existing global platforms are feature-rich but not always tailored for local payment flows and academic project constraints. CourShera addresses this by combining core MOOC features with payment verification support suitable for local contexts.

\section{Objectives}
The key objectives of CourShera are:
\begin{itemize}[leftmargin=*]
  \item Build a responsive frontend for course discovery, learning flow, and profile management.
  \item Build a backend API for authentication, course data retrieval, enrollment, payment, and user profile operations.
  \item Use a relational database with clear schema design and data relationships.
  \item Integrate third-party services where needed (Google OAuth, Supabase storage).
  \item Containerize services for reproducible setup and deployment readiness.
\end{itemize}

\section{Scope of the Implemented System}
The implemented system includes:
\begin{itemize}[leftmargin=*]
  \item User authentication: local registration/login and Google OAuth login.
  \item Course services: popular courses, recommendations, categories, category-wise listing, and individual course detail.
  \item Enrollment and payment flow: transaction verification and enrollment activation.
  \item User services: profile retrieval/update and saved card management.
  \item Quiz module pages: catalog, quiz home, taking quizzes, certificate, and past attempts.
\end{itemize}

\section{Updated Changes After Presentation Feedback}
The following report-level updates were made based on presentation feedback:
\begin{itemize}[leftmargin=*]
  \item All modeling chapters now include explicit sections for \textit{"Changes Made After Feedback"}.
  \item Screenshot placeholders were added for every diagram and mock UI chapter to ensure visual evidence in the final report.
  \item API documentation chapter now includes a dedicated subsection to list corrections between the documented specification and implemented endpoints.
  \item Feature Implementation chapter was expanded to include architecture, deployment, DevOps setup, and individual member contributions.
\end{itemize}

\clearpage

% =========================================================
\chapter{BPMN Diagram}

\textbf{Presented by:} [Member Name, Student ID]

\section{Final BPMN Diagram}
\textbf{Insert screenshot below:}
\begin{figure}[h!]
  \centering
  \includegraphics[width=0.95\textwidth]{assets/bpmn-final.png}
  \caption{Final BPMN Diagram}
\end{figure}

\section{Changes Made After Feedback}
\begin{itemize}[leftmargin=*]
  \item [Write what was changed in lanes, events, gateways, and error flows.]
\end{itemize}

\clearpage

% =========================================================
\chapter{Use Case Diagram}

\textbf{Presented by:} [Member Name, Student ID]

\section{Final Use Case Diagram}
\begin{figure}[h!]
  \centering
  \includegraphics[width=0.95\textwidth]{assets/use-case-final.png}
  \caption{Final Use Case Diagram}
\end{figure}

\section{Changes Made After Feedback}
\begin{itemize}[leftmargin=*]
  \item [Write actor/use-case relation updates and include/extend corrections.]
\end{itemize}

\clearpage

% =========================================================
\chapter{Class Diagram}

\textbf{Presented by:} [Member Name, Student ID]

\section{Final Class Diagram}
\begin{figure}[h!]
  \centering
  \includegraphics[width=0.95\textwidth]{assets/class-diagram-final.png}
  \caption{Final Class Diagram}
\end{figure}

\section{Changes Made After Feedback}
\begin{itemize}[leftmargin=*]
  \item [Write updates in class attributes, methods, cardinality, and associations.]
\end{itemize}

\clearpage

% =========================================================
\chapter{Mock UI}

\textbf{Presented by:} [Member Name, Student ID]

\section{Final Mock UI Screens}
\begin{figure}[h!]
  \centering
  \includegraphics[width=0.45\textwidth]{assets/mock-home-final.png}
  \includegraphics[width=0.45\textwidth]{assets/mock-course-final.png}
  \caption{Final Mock UI (Home and Course Details)}
\end{figure}

\section{Changes Made After Feedback}
\begin{itemize}[leftmargin=*]
  \item [Write updates made in navigation, hierarchy, consistency, and usability.] 
\end{itemize}

\clearpage

% =========================================================
\chapter{Collaboration Diagram}

\textbf{Presented by:} [Member Name, Student ID]

\section{Final Collaboration Diagram}
\begin{figure}[h!]
  \centering
  \includegraphics[width=0.95\textwidth]{assets/collaboration-final.png}
  \caption{Final Collaboration Diagram}
\end{figure}

\section{Changes Made After Feedback}
\begin{itemize}[leftmargin=*]
  \item [Write numbering/interaction corrections and object role clarifications.]
\end{itemize}

\clearpage

% =========================================================
\chapter{Sequence Diagram}

\textbf{Presented by:} [Member Name, Student ID]

\section{Final Sequence Diagram}
\begin{figure}[h!]
  \centering
  \includegraphics[width=0.95\textwidth]{assets/sequence-final.png}
  \caption{Final Sequence Diagram}
\end{figure}

\section{Changes Made After Feedback}
\begin{itemize}[leftmargin=*]
  \item [Write updates in lifelines, activation bars, alternate/exception flows.]
\end{itemize}

\clearpage

% =========================================================
\chapter{State Diagram}

\textbf{Presented by:} [Member Name, Student ID]

\section{Final State Diagram}
\begin{figure}[h!]
  \centering
  \includegraphics[width=0.95\textwidth]{assets/state-final.png}
  \caption{Final State Diagram}
\end{figure}

\section{Changes Made After Feedback}
\begin{itemize}[leftmargin=*]
  \item [Write updates to state transitions, guards, and terminal states.]
\end{itemize}

\clearpage

% =========================================================
\chapter{API Documentation}

\textbf{Presented by:} [Member Name, Student ID]

\section{Final API Documentation Snapshot}
\begin{figure}[h!]
  \centering
  \includegraphics[width=0.95\textwidth]{assets/api-doc-final.png}
  \caption{Final API Documentation Screenshot}
\end{figure}

\section{Changes Made After Feedback}
\begin{itemize}[leftmargin=*]
  \item [Document corrected base URL/environment details.] 
  \item [Document endpoint path updates and request/response schema fixes.] 
  \item [Document authentication and error-response clarifications.] 
\end{itemize}

\section{Reference File}
OpenAPI specification file used in this project: \texttt{api\_docs.yaml}

\clearpage

% =========================================================
\chapter{Feature Implementation}

\textbf{Presented by:} [Member Name, Student ID]

\section{Feature Name and Description}
\textbf{Feature Name:} Integrated Course Discovery, Enrollment, and Learner Profile Management in CourShera.

\textbf{Description:}
This feature combines the main user journey of the system:
\begin{itemize}[leftmargin=*]
  \item Discovering courses using popular lists, recommendations, categories, and detailed course pages.
  \item Authenticating users through local credentials or Google OAuth.
  \item Completing enrollment through verified payment transaction matching.
  \item Managing user profile data and saved card records.
  \item Accessing quiz-center flows to support learning engagement.
\end{itemize}

\section{GitHub Repository Link(s)}
\begin{itemize}[leftmargin=*]
  \item Main repository: \url{[Add your GitHub repository URL here]}
\end{itemize}

\textit{Note: No git remote URL is currently configured in this local workspace, so the repository URL must be inserted manually.}

\section{Hosted Website Link}
\begin{itemize}[leftmargin=*]
  \item Frontend URL: \url{[Add deployed frontend URL or write "Not deployed"]}
  \item Backend API URL: \url{[Add deployed backend URL or write "Not deployed"]}
\end{itemize}

\section{Screenshots of Different Pages}
\begin{figure}[h!]
  \centering
  \includegraphics[width=0.48\textwidth]{assets/ss-home.png}
  \includegraphics[width=0.48\textwidth]{assets/ss-category.png}
  \caption{Home Page and Category Page}
\end{figure}

\begin{figure}[h!]
  \centering
  \includegraphics[width=0.48\textwidth]{assets/ss-course-outline.png}
  \includegraphics[width=0.48\textwidth]{assets/ss-checkout.png}
  \caption{Course Outline and Checkout Page}
\end{figure}

\begin{figure}[h!]
  \centering
  \includegraphics[width=0.48\textwidth]{assets/ss-profile.png}
  \includegraphics[width=0.48\textwidth]{assets/ss-quiz-catalog.png}
  \caption{Profile Page and Quiz Catalog Page}
\end{figure}

\section{System Architecture Overview}
\subsection{Overall Architecture Used}
The implementation follows a \textbf{layered monolithic architecture} with clear frontend-backend-database separation:
\begin{itemize}[leftmargin=*]
  \item \textbf{Presentation Layer:} React (Vite) single-page frontend.
  \item \textbf{Application/Service Layer:} Express.js backend exposing REST endpoints.
  \item \textbf{Data Access Layer:} Prisma ORM for PostgreSQL operations.
  \item \textbf{Data Layer:} PostgreSQL database schema with entities such as clients, courses, enrollments, payments, reviews, and categories.
\end{itemize}

\subsection{Technologies and Frameworks}
\begin{itemize}[leftmargin=*]
  \item \textbf{Frontend:} React 19, React Router, Vite, CSS modules/styles.
  \item \textbf{Backend:} Node.js, Express 5, Passport (Google OAuth + local strategy), cookie-session.
  \item \textbf{Database:} PostgreSQL with Prisma ORM.
  \item \textbf{Storage/External Service:} Supabase storage for user avatar uploads.
  \item \textbf{Security/Utility:} bcrypt for password hashing, JSON/session-based auth flows.
\end{itemize}

\subsection{High-Level Runtime Flow}
\begin{enumerate}[leftmargin=*]
  \item User interacts with React frontend routes (home, category, course, checkout, profile, quiz pages).
  \item Frontend calls backend REST APIs (e.g., /auth, /courses, /payment, /me).
  \item Backend validates session/auth, executes business logic, and queries PostgreSQL through Prisma.
  \item Response data is returned to frontend for rendering and user interaction.
\end{enumerate}

\section{Deployment Details}
\subsection{Where Components Are Hosted}
Current implementation is configured for \textbf{local and containerized deployment}:
\begin{itemize}[leftmargin=*]
  \item Frontend container maps to port 3000 (Nginx serving built app).
  \item Backend container maps to port 5000.
  \item Database uses PostgreSQL connection configured via environment variables (project setup indicates Supabase-hosted PostgreSQL in use for development).
\end{itemize}

\subsection{Cloud Platform Used}
\begin{itemize}[leftmargin=*]
  \item Database/Storage: Supabase services are integrated in the codebase.
  \item Frontend/Backend public cloud deployment URL: \textit{Not finalized in this repository snapshot; add production/staging links here if available.}
\end{itemize}

\subsection{DevOps / Deployment Setup}
\begin{itemize}[leftmargin=*]
  \item \textbf{Containerization:} Dockerfiles exist for frontend and backend.
  \item \textbf{Orchestration:} docker-compose configuration provided to run full stack.
  \item \textbf{CI/CD:} No CI/CD pipeline configuration was found in this repository snapshot. Add GitHub Actions/GitLab CI details here if configured later.
\end{itemize}

\section{Contribution of Each Member}
\begin{longtable}{|p{0.24\textwidth}|p{0.70\textwidth}|}
\hline
\textbf{Member (Name, ID)} & \textbf{Contribution in Feature Implementation} \\
\hline
[Member 1] & [Frontend module(s), pages/components, UI integration] \\
\hline
[Member 2] & [Backend routes/services, authentication/payment logic] \\
\hline
[Member 3] & [Database design, Prisma schema, API docs] \\
\hline
[Member 4] & [Testing, deployment, integration, documentation] \\
\hline
\end{longtable}

\section{Implemented Endpoints and Modules (Summary)}
\begin{itemize}[leftmargin=*]
  \item \textbf{Auth Module:} local register/login, Google OAuth login success/failure/logout.
  \item \textbf{Courses Module:} categories, category courses, in-progress courses, recommendations, popular courses, single-course detail.
  \item \textbf{Payment Module:} transaction verification and enrollment activation.
  \item \textbf{Profile Module:} profile fetch/update, saved cards add/list/delete.
\end{itemize}

\section{Changes Made After Presentation Feedback}
\begin{itemize}[leftmargin=*]
  \item Expanded this chapter to include architecture style, stack details, deployment notes, and DevOps setup.
  \item Added explicit placeholders for repository links, hosted links, and page screenshots.
  \item Added a structured contribution table to clearly show each member's implementation role.
\end{itemize}

\end{document}
